Let $f\colon X \to X$ be a self map of a compact CW-complex with induced pullback map
$H^i(f)\colon H^i(X;\QQ)\to H^i(X;\QQ)$ on (rational) singular cohomology. Since
$X$ is compact, each of the cohomology groups is finitely generated and there but finitely
many non-zero groups, so one defines the \emph{Lefschetz number} of $f$ to be the alernating sum:
\[ L(f)\coloneq\sum_{i\geq 0} \mathrm{tr}\left(H^i(f)\right)\]
The celebrated \emph{Lefschetz fixed point theorem} states that if the set of fixed points
$X^f\subseteq X$ is finite and discrete, there exists some naturally defined integers 
$\lambda^f(x)$ for $x\in X^f$ such that:
\[  \sum_{x \in X^f} \lambda^f(x) = L(f)\,.\]
The allure of this theorem is straightforward; the right hand side only depends
on the homotopy type of $f$, why the left hand side varies wildly (and indeed may not even 
be defined) across different but homotopic maps.

 The goal of this note is to prove the following far reaching generalization:
\begin{thm}\label{thm: main}
  Let $p\colon X \to \base$ be a proper shape submersion of locally compact Hausdorff spaces
  and let $f\colon X \to X$ be a self map of $X$ over $\base$. Write $X^f\subseteq X$ for
  the fixed points and $\cL^f X$ for the homotopy fixed points\footnote{i.e.~the space of pairs $(x,\gamma)$ where $x\in X$
  and $\gamma$ is a path from $x$ to $f(x)$.} of $f$.
  Moreover, let
  \[\mdef{L(f)} \in H^0_{\mathrm{c}}(\base)\]
  denote trace of the morphism induced by $f$ on the relative cohomology $\cptSec{X}[\base]$.
  There exists a natural\footnote{Elaborate} commutative diagram in the $\infty$-category
  $\Shv(\base,\Sp)$ of the form:
  \[
\begin{tikzcd}
	{\cptSec{X^f}[\base]^\vee} & {\SS_{\base}[\cL^f X]} \\
	{\SS_{\base}} & {\SS_{\base}}
	\arrow["{\ass_f}", from=1-1, to=1-2]
	\arrow["\pi", from=1-2, to=2-2]
	\arrow["{\lambda_f}", from=2-1, to=1-1]
	\arrow["{\rho_f}"{description}, from=2-1, to=1-2]
	\arrow["{L(f)}"', from=2-1, to=2-2]
\end{tikzcd}
\]
where $\pi$ is the pushforward along the map $\cL^fX \to \base$ and $\ass_f$ is induced
by the map taking any $x\in X^f$ to the pair $(x,\mathrm{const}_x).$
\end{thm}

\begin{rmk}
  We make some observations about this theorem. 
  \begin{enumerate}
    \item Note that we make \emph{no assumptions} on the set of fixed points $X^f$ and its topology.
      In particular, we deduce the weak Lefschetz theorem, namely that $L(f)\neq 0$ implies
      that any map homotopic to $f$ has a fixed point in this generality.
    \item In the case $M=\pt$, the factorization of $L(f)$ as composite 
      \[ \SS \xrightarrow{\lambda_f} \Gamma(X^f) \to \SS\]
      readily recovers the classical Lefschetz fixed point theorem whenever $X^f$ is discrete.
    \item The map $\rho_f$ fiberwise agrees with the \emph{Reidemeister trace}. 
    \item If $M=\pt$ and $X^f\subseteq X$ is finite and discrete, we further deduce the 
      Lefschetz--Nielsen theorem, providing a lower bound for the number of fixed points
      of any map homotopic to $f$ in terms of the Reidemeister trace.
  \end{enumerate}
\end{rmk}
